% =====================================================================
%  Stage 1, раздел 3.
%  Обзор подходов к моделированию атак, генеративных моделей
%  и обоснование архитектуры ИИ-злоумышленника.
% =====================================================================

\section{Обзор подходов к моделированию атак и обоснование архитектуры ИИ-злоумышленника}

\subsection{Назначение ИИ-злоумышленника в системе оценки NIDS}

Под ИИ-злоумышленником (далее --- агентом) в работе понимается активный
программный компонент, формирующий поток сетевых событий с управляемым
переключением режимов \emph{норма / атака / дрейф} и подающий этот поток на
вход детектора. Назначение агента --- обеспечить воспроизводимое и
параметризованное окружение для оценки качества и устойчивости детектора в
условиях, приближенных к эксплуатационным.

В отличие от пассивных датасетов, агент позволяет:
\begin{itemize}
    \item получить контролируемую \emph{ground truth} в потоковом сценарии,
    включая моменты переключения режимов;
    \item варьировать частоты атак, длительность сессий и параметры сетевых
    шаблонов в широких пределах;
    \item моделировать дрейф распределений, проверяя устойчивость детектора
    и работу drift-мониторинга;
    \item обеспечить независимость тестового сценария от датасета,
    использованного для обучения, что снижает риск переобучения на специфике
    UNSW-NB15.
\end{itemize}

Отдельно подчеркнём, что речь идёт о \emph{симуляционном агенте на уровне
flow-признаков}: целью не является эмуляция реальных атак на работающую
инфраструктуру. Это разделение существенно с точки зрения этики и
юридического статуса исследования.

\subsection{Подходы к моделированию атак: систематика}

В литературе по ML-NIDS выделяются три устойчивых семейства подходов к
формированию атакующего трафика для тестирования детекторов.

\begin{enumerate}
    \item \textbf{Шаблонная (\eng{rule-based, template-based}) генерация.}
    Атакующие сценарии задаются параметрическими шаблонами: распределения
    межпакетных интервалов, длины потоков, наборы флагов TCP, скорости
    передачи. Применение --- широкое: от классических работ DARPA/KDD до
    инструментов \eng{traffic replay} вроде Tcpreplay и
    Hping3~\cite{lippmann2000darpa, gharib2016evalframework}.
    \item \textbf{Генеративные модели.} Использование GAN, cGAN, WGAN,
    VAE и автоэнкодеров для синтеза реалистичных flow-векторов и
    последовательностей. Эталонная работа --- Ring и
    соавт.~\cite{ring2019flowgan}; вариационные подходы для IoT-сетей
    предложены в~\cite{lopez2017cvae}.
    \item \textbf{Adversarial-агенты на основе обучения с подкреплением.}
    Агент учится формировать последовательности действий, обходящих
    детектор, в форме, аналогичной white-/grey-box атаке на классификатор;
    эталонная работа --- Apruzzese и соавт.~\cite{apruzzese2020adversarial}.
\end{enumerate}

В \tabref{tab:attacker-approaches} приведено сводное теоретическое сравнение
этих семейств по пяти эксплуатационно значимым критериям.

\begin{table}[H]
\centering
\caption{Сравнение подходов к моделированию атак для оценки NIDS}
\label{tab:attacker-approaches}
\renewcommand{\arraystretch}{1.2}
\begin{tabularx}{\linewidth}{@{}lXXX@{}}
\toprule
\textbf{Критерий} & \textbf{Шаблонная} & \textbf{Генеративная (GAN/VAE)} &
\textbf{Adversarial-RL} \\
\midrule
Воспроизводимость       & высокая  & средняя; зависит от стабильности обучения &
низкая; зависит от инициализации политики \\
\addlinespace
Семантическая
интерпретируемость      & прямая (параметры~$\to$~атака) & косвенная (через
обученное распределение) & низкая (политика — «чёрный ящик») \\
\addlinespace
Вариативность           & ограничена шаблонами; усиливается рандомизацией параметров &
высокая; ограничена режимом коллапса GAN & высокая, но смещена в сторону
обхода конкретного детектора \\
\addlinespace
Стоимость разработки    & низкая & высокая (стабильное обучение GAN) &
очень высокая (среда + reward design + обучение) \\
\addlinespace
Совместимость с этикой
и охватом ВКР           & полная & требует валидации реалистичности & требует
явного контроля; возможна оптимизация под обход детектора \\
\bottomrule
\end{tabularx}
\end{table}

\subsection{Шаблонная генерация: математическая модель}

Шаблон атаки описывается параметрическим вектором
\(\theta \in \Theta\), задающим распределения базовых характеристик
flow-сессий: длительность \(d \sim F_d(\cdot;\theta)\), скорость пакетов
\(\lambda \sim F_\lambda(\cdot;\theta)\), доля флагов SYN/ACK и др. Поток
наблюдений на интервале \([0, T]\) задаётся точечным процессом
\(\mathcal{N}(t;\theta)\), параметры которого определяются типом атаки.
Например, для DDoS volumetric SYN-flood:
\begin{equation}
\lambda(t) = \lambda_0 \cdot \mathbb{I}[t \in [t_a, t_a + \Delta]] + \lambda_n,
\label{eq:ddos-rate}
\end{equation}
где \(\lambda_n\) --- интенсивность фонового потока, \(\lambda_0\) ---
интенсивность атаки, \([t_a, t_a + \Delta]\) --- временной интервал атаки.
Шаблон позволяет изменять \(\lambda_0\), \(\Delta\) и долю \emph{SYN}-флагов,
получая семейство сценариев одной и той же атаки.

Аналогичные шаблоны определены для остальных классов: Reconnaissance задаётся
заметанием destination-портов с фиксированной интенсивностью обращений к
одному источнику; Brute Force --- регулярными подключениями к одному порту с
малым числом байт за сессию; Exfiltration --- длинной устойчивой передачей с
дисбалансом \(b_{out} \gg b_{in}\); Internal Reconnaissance --- широким
фан-аутом потоков от одного источника к множеству destinations с типичной
поведенческой подписью.

Преимущество шаблонного подхода --- прозрачность ground truth: момент атаки,
её тип и параметры однозначно зафиксированы конфигурацией. Ограничение ---
ограниченная вариативность; компенсируется введением рандомизации параметров.

\subsection{Генеративные модели для синтеза трафика: теоретическая база}

\subsubsection{GAN: формализм}

Классическая постановка GAN~\cite{goodfellow2014gan}:
\begin{equation}
\min_{G}\max_{D}
V(D,G) = \mathbb{E}_{x \sim p_{\text{data}}}[\log D(x)] +
         \mathbb{E}_{z \sim p_z}[\log(1 - D(G(z)))],
\label{eq:gan}
\end{equation}
где \(G\colon \mathcal{Z} \to \mathcal{X}\) --- генератор,
\(D\colon \mathcal{X} \to [0,1]\) --- дискриминатор,
\(p_z\) --- априорное распределение в латентном пространстве.
Минимакс \eqnref{eq:gan} в идеале сводится к равенству распределений
\(G(z) \sim p_{\text{data}}\) при оптимальном дискриминаторе.

\subsubsection{Conditional GAN}

Conditional GAN~\cite{mirza2014cgan} вводит обусловливающую переменную \(c\)
(тип атаки):
\begin{equation}
\min_{G}\max_{D}
V(D,G) = \mathbb{E}_{x,c}[\log D(x,c)] +
         \mathbb{E}_{z,c}[\log(1 - D(G(z,c), c))].
\label{eq:cgan}
\end{equation}
Это позволяет генерировать целенаправленные атакующие сценарии конкретного
класса, что и составляет содержание подхода
Ring и соавт.~\cite{ring2019flowgan}.

\subsubsection{WGAN-GP}

Обучение GAN страдает от нестабильности и режима коллапса (\eng{mode
collapse}). Wasserstein GAN~\cite{arjovsky2017wgan} заменяет JS-расстояние
расстоянием Earth Mover (Wasserstein-1):
\begin{equation}
W(p, q) = \inf_{\gamma \in \Pi(p, q)} \mathbb{E}_{(x,y) \sim \gamma}\|x - y\|.
\label{eq:wgan}
\end{equation}
Расширение WGAN-GP~\cite{gulrajani2017wgangp} вводит градиентный штраф,
обеспечивающий 1-липшицевость дискриминатора:
\begin{equation}
\mathcal{L}_{\text{GP}} = \lambda\,
\mathbb{E}_{\hat{x}}\bigl[(\|\nabla_{\hat{x}} D(\hat{x})\|_2 - 1)^2\bigr].
\label{eq:wgangp}
\end{equation}
WGAN-GP --- стандартный выбор для синтеза табличных flow-данных.

\subsubsection{Применение GAN-семейства в NIDS}

Систематический обзор Ring и соавт.~\cite{ring2019flowgan} показал, что
GAN способен синтезировать реалистичные flow-векторы с распределениями,
неотличимыми по ряду статистических критериев. Однако на практике
зафиксированы:

\begin{itemize}
    \item Чувствительность качества к гиперпараметрам и инициализации.
    \item Режим коллапса для редких классов (Worms, Backdoors), где
    обучающих данных мало.
    \item Сложность валидации семантической согласованности
    сгенерированного трафика (например, согласованности TCP-флагов и
    длительности).
\end{itemize}

Эти ограничения существенны в условиях ВКР, где требуется обеспечить
воспроизводимое тестирование и однозначную ground truth.

\subsection{Adversarial-агенты на основе обучения с подкреплением}

Apruzzese и соавт.~\cite{apruzzese2020adversarial} рассматривают атакующего
как агента в среде, где \(\text{state}\) --- текущее состояние
взаимодействия с детектором, \(\text{action}\) --- модификация параметров
исходящего трафика, \(\text{reward}\) --- успех обхода детектора.
Стандартная политика обучается алгоритмами PPO, DQN или DDPG.

С точки зрения настоящей работы такой подход:
\begin{itemize}
    \item даёт мощный инструмент для аудита устойчивости детектора;
    \item трудоёмок в реализации (требует среды, симулирующей ответ
    детектора в режиме обучения);
    \item смещает поведение агента в сторону обхода конкретного
    детектора, что не совпадает с задачей формирования \emph{внешних}
    тестовых сценариев;
    \item требует расширенного этического контроля и явного отказа от
    обобщения на реальные системы.
\end{itemize}

В рамках ВКР этот подход рассматривается как направление дальнейшего
развития, но не как реализуемый компонент.

\subsection{Concept drift: моделирование и оценка}

Феномен \emph{concept drift} --- изменение распределений данных и/или
условной зависимости меток от признаков во времени --- является основным
источником деградации NIDS в эксплуатации~\cite{gama2014drift,
pendlebury2019tesseract, andresini2021insomnia}.

Различают четыре типа дрейфа~\cite{gama2014drift}:
\begin{enumerate}
    \item \textbf{Sudden drift} --- резкое изменение распределения
    \(p(x, y)\);
    \item \textbf{Incremental drift} --- постепенное смещение во времени;
    \item \textbf{Gradual drift} --- частая смена режимов между двумя
    распределениями;
    \item \textbf{Recurring drift} --- циклически возвращающиеся режимы.
\end{enumerate}

Для целей настоящей работы дрейф моделируется как преобразование
\(\mathcal{T}\colon X \to X\), применяемое к подмножеству признаков. Три
варианта инжектора реализуются в Этапе 3 и используются в Эксперименте 5
Этапа 4:

\begin{itemize}
    \item \textbf{Covariate shift}:
    \(x' = \mathcal{T}_{\text{cov}}(x; \mu, \sigma) = \mu + \sigma \odot x\),
    что соответствует изменению масштаба и положения признаков без
    изменения метки;
    \item \textbf{Prior shift}: изменение долей классов в потоке
    через политику переключения режимов \emph{ATTACK(type)} в
    конечном автомате агента;
    \item \textbf{Concept drift}: изменение зависимости \(p(y \mid x)\)
    через модификацию параметров шаблонов атак (например, понижение
    \(\lambda_0\) DDoS до уровней, граничащих с нормой).
\end{itemize}

Контроль интенсивности дрейфа осуществляется через индексы
\(\PSI\)~\eqnref{eq:psi} и \(\MMD\)~\eqnref{eq:mmd}, что обеспечивает
сопоставимость экспериментов.

\subsection{Архитектура ИИ-злоумышленника, принимаемая в работе}

С учётом проведённого сравнения и условий, формируемых задачей ВКР, в работе
принимается \textbf{гибридная} архитектура агента:

\[
\boxed{\;\text{Templates} \;+\; \text{FSM-планировщик} \;+\; \text{Drift-инжектор}\;}
\]

\noindent без использования генеративных моделей в качестве реализуемого
компонента. Генеративные подходы (GAN/cGAN/WGAN-GP) рассматриваются как
теоретическая альтернатива, обоснованно отвергаемая по следующим причинам.

\begin{enumerate}
    \item \textbf{Воспроизводимость.} Шаблонный агент даёт однозначную
    \emph{ground truth} (момент, тип, параметры атаки). GAN-агент, напротив,
    требует фиксации seeds, состояния обучения и контроля режима коллапса,
    что существенно усложняет научную верификацию результатов.
    \item \textbf{Стабильность обучения.} В условиях ВКР с ограниченным числом
    атакующих образцов (особенно для редких классов в UNSW-NB15) GAN-обучение
    подвержено режиму коллапса, что приводит к нерепрезентативным сценариям и
    систематическим ошибкам экспериментальной части.
    \item \textbf{Семантическая согласованность.} Шаблоны явно поддерживают
    согласованность между признаками (например, TCP-флагами и
    длительностью), тогда как GAN-генерация на табличных flow-данных требует
    отдельных верификационных проверок (TSTR, маргиналы, корреляции).
    \item \textbf{Привязка к таксономии.} Шаблоны позволяют детерминированно
    соответствовать классам UNSW-NB15, что упрощает интерпретацию ошибок
    детектора и анализ FP/FN.
\end{enumerate}

При этом сохраняемое в работе теоретическое описание GAN-семейства
обеспечивает полноту обзора и явно фиксирует направление дальнейшего
развития методики --- расширение шаблонной библиотеки за счёт обученной
cGAN-составляющей.

\subsection{Конечно-автоматная политика агента}

Состояния планировщика образуют конечный автомат:
\[
\mathcal{S} = \{\text{NORMAL}\} \cup
\{\text{ATTACK}(c)\}_{c \in \mathcal{C}} \cup
\{\text{DRIFT}(k)\}_{k \in \mathcal{K}},
\]
где \(\mathcal{C}\) --- множество типов атак из таксономии UNSW-NB15
(в реализации --- подмножество с реалистичной моделью flow-проявлений),
\(\mathcal{K}\) --- множество типов дрейфа (covariate / prior / concept).

Переходы между состояниями определяются дискретной политикой
\(\pi\colon \mathcal{S} \to \Delta(\mathcal{S})\), параметризованной
\emph{вероятностями переключения} и \emph{минимальными длительностями
сессий} для каждого состояния. Такая параметризация обеспечивает:
\begin{itemize}
    \item возможность задавать долю атакующего трафика в потоке (\emph{attack
    ratio});
    \item возможность задавать частоту смены режимов (\emph{flicker rate});
    \item полную воспроизводимость прогона при фиксированном seed.
\end{itemize}

Формально политика для перехода из NORMAL имеет вид
\[
\pi(\text{NORMAL}) =
\bigl((1-p_a-p_d)\!\cdot\!\text{NORMAL}\bigr)
\oplus \sum_{c}\!p_a^{(c)}\!\cdot\!\text{ATTACK}(c)
\oplus \sum_{k}\!p_d^{(k)}\!\cdot\!\text{DRIFT}(k),
\]
где \(\oplus\) --- смесь распределений. Минимальные длительности
\(d_{\min}^{s}\) для каждого состояния гарантируют отсутствие
``мерцающих'' переключений, нерепрезентативных в эксплуатации.

\subsection{Ground truth, метаданные и протокол валидации агента}

Каждое сгенерированное окно сопровождается записью \(\langle t, s, \theta\rangle\),
где \(t\) --- метка времени, \(s\) --- состояние FSM, \(\theta\) ---
параметры шаблона. Этот журнал служит неизменной \emph{ground truth} для
вычисления метрик (особенно time-to-detect и FP-per-time).

Протокол валидации агента включает:
\begin{enumerate}
    \item \textbf{Маргинальное соответствие.} Сравнение распределений ключевых
    признаков (\emph{duration, sbytes, dbytes, sttl, dur}) в режиме NORMAL
    с распределениями нормального трафика UNSW-NB15 по тестам
    Колмогорова--Смирнова и MMD. Цель --- подтвердить, что фон агента не
    тривиально отличается от референсной нормы.
    \item \textbf{Семантическую согласованность атак.} Проверка соответствия
    флагов TCP, длительности и интенсивности шаблонных атак данным UNSW-NB15
    для соответствующего класса.
    \item \textbf{Контроль интенсивности дрейфа.} Расчёт PSI/MMD между
    эталонным и текущим окном, попадание в заданные диапазоны при
    включении DRIFT-состояний.
\end{enumerate}

\subsection{Этические и юридические аспекты}

Все процедуры моделирования атак выполняются исключительно в составе
автономного программного стенда; целью не является эмуляция реальных атак
на работающие сети или сервисы. Сценарии формируются в виде синтетических
flow-векторов, не передаваемых в реальную сетевую среду. В работе
отсутствуют операции, направленные на обход эксплуатируемых средств защиты
третьих лиц, и отсутствуют рекомендации по практическому применению
атакующих техник. Это соответствует общепринятой в академическом сообществе
практике исследований adversarial-NIDS~\cite{apruzzese2020adversarial,
ring2019flowgan}.

\subsection{Промежуточные выводы по разделу}

\begin{enumerate}
    \item Систематизированы три семейства подходов к моделированию атак
    для оценки NIDS: шаблонные, генеративные и adversarial-агенты на
    обучении с подкреплением. Каждое семейство охарактеризовано по пяти
    критериям; отмечено, что шаблонный подход обеспечивает наивысшую
    воспроизводимость и интерпретируемость, но ограничен по вариативности.
    \item Описана математическая база генеративного семейства (GAN, cGAN,
    WGAN-GP), приведены целевые функции и зафиксированы ограничения
    применимости в условиях ВКР: нестабильность обучения, режим коллапса
    для редких классов, сложность валидации семантической согласованности.
    \item Обоснован выбор гибридной архитектуры агента: \emph{шаблоны +
    конечный автомат + drift-инжектор} без cGAN. Описаны математические
    модели шаблонов основных классов атак (DDoS, Reconnaissance, Brute Force,
    Exfiltration, Internal Reconnaissance) и формальная политика FSM.
    \item Зафиксированы три типа моделируемого дрейфа (covariate, prior,
    concept) и метрики его контроля (PSI, MMD), что обеспечивает связность
    с протоколом экспериментов главы~4.
    \item Сформулирован протокол валидации агента (маргинальное соответствие,
    семантическая согласованность, контроль интенсивности дрейфа) и
    зафиксированы этические рамки исследования.
\end{enumerate}
